% Drawing 78/07-4 -- folder 78, batch 07, archive page 129 (PDF page 130).
% Transcription: deux arcs se terminant chacun par un point épais, au-dessus d'un segment horizontal marqué à droite d'un trait vertical et de la lettre $H$ (un astérisque raturé à gauche).
% Kind: two arcs over a segment. Method 2 (schematic redraw from the facsimile). Status: ok.
% Notes: Two arcs, each ending in a heavy point (the left one drawn as two strokes from its point), above a horizontal segment marked on the right by a vertical stroke and H; a struck asterisk on the left (drawn grey).
% Coordinates: 1 unit = 1 mm = 1/100 of the crop width (bbox in page fractions: [0.3, 0.36, 0.57, 0.46]).
% Redrawn by Claude (claude-opus-5-5), issue #73, 2026-11-07; not checked by a human; not compiled (no TeX in the session).
% Rights: reproduces part of the work -- behind RENDER_TRANSCRIPTS, never published by itself.
\documentclass[tikz,border=3pt]{standalone}
\usepackage{amsmath,amssymb}
\usetikzlibrary{patterns,arrows.meta}
\tikzset{pencil/.style={black!45}, struck/.style={black!35}, mark/.style={circle,draw,inner sep=1.3pt}, pt/.style={circle,fill,inner sep=0.9pt}, opt/.style={circle,draw,fill=white,inner sep=1.1pt}}
\begin{document}
\begin{tikzpicture}[x=1mm,y=1mm,line cap=round,line join=round,font=\scriptsize]
\node[inner sep=1pt,struck] at (13.7,-7.5) {\small $\ast$};
\draw plot[smooth, tension=0.55] coordinates {(19.3,-26.4) (27,-17.5) (34.8,-10.8)};
\draw (19.3,-26.4) -- (41.5,-16.4);
\node[pt] at (19.3,-26.4) {};
\draw plot[smooth, tension=0.55] coordinates {(54.8,-11.9) (61.5,-19.7) (72.6,-24.1)};
\draw (67,-9.7) -- (72.6,-24.1);
\node[pt] at (72.6,-24.1) {};
\draw (10.4,-44.1) -- (70.4,-45.2);
\draw (79.3,-40.8) -- (80.4,-48.6);
\node[inner sep=1pt] at (88.1,-44.1) {\tiny $H$};
\end{tikzpicture}
\end{document}
